\documentclass[a4paper,oneside,11pt]{book}

% Package yang diperlukan
\usepackage{graphicx}
\usepackage{titling}
\usepackage{geometry}
\usepackage{listings}
\usepackage{xcolor}
\usepackage{float}
\usepackage{hyperref}
\usepackage{indentfirst}

% Pengaturan margin
\geometry{
    left=3cm,
    right=2.5cm,
    top=3cm,
    bottom=3cm
}

% Pengaturan untuk kode HTML
\definecolor{codegreen}{rgb}{0,0.6,0}
\definecolor{codegray}{rgb}{0.5,0.5,0.5}
\definecolor{codepurple}{rgb}{0.58,0,0.82}
\definecolor{backcolour}{rgb}{0.95,0.95,0.92}

\lstdefinestyle{htmlstyle}{
    backgroundcolor=\color{backcolour},   
    commentstyle=\color{codegreen},
    keywordstyle=\color{codepurple},
    numberstyle=\tiny\color{codegray},
    stringstyle=\color{codegreen},
    basicstyle=\ttfamily\footnotesize,
    breakatwhitespace=false,         
    breaklines=true,                 
    captionpos=b,                    
    keepspaces=true,                 
    numbers=left,                    
    numbersep=5pt,                  
    showspaces=false,                
    showstringspaces=false,
    showtabs=false,                  
    tabsize=2,
    frame=single
}

\title{Laporan Praktikum \\ Praktikum Pemrograman Web 1\\
Pertemuan 2\\
Dasar-dasar HTML}
\author{Nama Mahasiswa: xxxxx(NIM)}
\date{\today}

\begin{document}

% Halaman judul
\begin{titlingpage}
\begin{center}
\vspace{4cm}
\begin{huge} 
\textbf{\thetitle} \\
\end{huge}
\vspace{2cm}
\includegraphics[height=8cm]{lambang ugm.png}\\
\begin{Large}
\vspace{4cm}
\theauthor\\
Kelas A1 \\
\end{Large}
\begin{large}
Dosen Pengampu: Dinar Nugroho Pratomo, S.Kom., M.IM., M.Cs.\\
\end{large}
\vspace{2cm}
\thedate
\end{center}
\end{titlingpage}

% Daftar Isi
\tableofcontents
\newpage

\chapter{Tujuan Praktikum}
\begin{enumerate}
    \item Mahasiswa mampu memahami.....
\end{enumerate}

\chapter{Dasar Teori}
\section{1}

\section{2}
\begin{lstlisting}[language=HTML, style=htmlstyle]
\end{lstlisting}

Seperti yang terlihat di gambar \ref{pic:komputer} ini adalah gambar komputer
\begin{figure}[H]
    \centering
    \includegraphics[height=2cm]{komputer.jpg}
    \caption{Gambar Komputer}
    \label{pic:komputer}
\end{figure}

\chapter{}
\section{3}

\chapter{Kesimpulan}
[Tuliskan kesimpulan yang diperoleh dari praktikum ini, sesuaikan dengan tujuan praktikum]

\chapter{Daftar Pustaka}
\begin{enumerate}
    
\end{enumerate}

\end{document}